\uuid{EfwX}
\exo7id{1606}
\titre{exo7 1606}
\auteur{barraud}
\organisation{exo7}
\datecreate{2003-09-01}
\isIndication{false}
\isCorrection{false}
\chapitre{Réduction d'endomorphisme, polynôme annulateur}
\sousChapitre{Valeur propre, vecteur propre}
\cpge{mpsi-mp2i.calcul-matriciel-systemes.e, pcsi.determinants.c, ptsi.determinants.a, mpsi-mp2i.determinants.d, pcsi.calcul-matriciel-systemes-lineaires.d, ptsi.calcul-matriciel-systemes-lineaires.d, bcpst1.matrices, mp-mpi.reduction-endomorphismes.b, pc.reduction.a, psi.reduction-endomorphismes.a, pt.reduction-endomorphismes.a, mp-mpi.reduction-endomorphismes.c, pc.reduction.b, psi.reduction-endomorphismes.b, pt.reduction-endomorphismes.b, bcpst2.valeurs-propres-vecteurs-propres.a, bcpst2.valeurs-propres-vecteurs-propres.b}
\module{Algèbre}
\niveau{L2}
\difficulte{}

\contenu{
\texte{
Soit $J$ la matrice $$J=
\begin{pmatrix}
  1 &\cdots &1\\
  \vdots&   &\vdots\\
  1 &\cdots &1
\end{pmatrix}
$$.
}
\begin{enumerate}
    \item \question{Trouver une relation entre $J$ et  $J^{2}$.}
    \item \question{En déduire les valeurs propres de $J$ et calculer leurs multiplicités.}
    \item \question{Donner le polynôme caractéristique de $J$.}
\end{enumerate}
}
